% !TEX program = xelatex
% ============================================================
%  Arabic Academic Article Template — Nibras Center
%  Compile with:
%    xelatex main.tex
%    bibtex main
%    xelatex main.tex
%    xelatex main.tex
%  Based on the BCI_Academic_arabic_Book + ASPU_Proposal_AR setups.
% ============================================================
%  Features:
%    - A4, 11pt, article class
%    - Title block with author, affiliation, abstract
%    - Glossaries + acronyms (\gls{}, \acrfull{}, \acrshort{})
%    - BibTeX bibliography (apalike style)
%    - Sections/subsections/subsubsections
%    - Fancy headers with title + author
% ============================================================

\documentclass[11pt,a4paper]{article}

% ---- Geometry ----
\usepackage[a4paper,margin=2.5cm,top=2.5cm,bottom=2.5cm,headheight=15pt]{geometry}

% ---- Core packages (order matters for bidi) ----
\usepackage{url}
\usepackage{amsmath}
\usepackage{graphicx}
\usepackage{array}
\usepackage{booktabs}
\usepackage{tabularx}
\usepackage{multirow}
\usepackage{enumitem}
\usepackage{float}
\usepackage{caption}
\usepackage{titlesec}
\usepackage{fancyhdr}
\usepackage{setspace}
\usepackage{supertabular}
\usepackage{longtable}

% ---- RTL / Arabic language support (must load before hyperref) ----
\usepackage{bidi}
\usepackage[no-math]{fontspec}
\usepackage[quiet]{polyglossia}

% ---- Fonts ----
\setmainfont[Scale=1]{Amiri-Regular.ttf}
\newfontfamily\arabicfont[Script=Arabic,Scale=0.95]{Amiri-Regular.ttf}
\newfontfamily\englishfont[Scale=0.95]{TeX Gyre Termes}
\setmonofont{Amiri-Regular.ttf}

% ---- Languages ----
\setdefaultlanguage[calendar=gregorian,hijricorrection=1]{arabic}
\setotherlanguage[variant=british]{english}

% ---- Hyperref (load after polyglossia/bidi) ----
\usepackage[breaklinks,hidelinks]{hyperref}
\hypersetup{
  pdftitle={عنوان المقال},
  pdfauthor={د. اسم المؤلف}
}

% ---- Glossaries + Acronyms ----
\usepackage[acronym,toc]{glossaries}
\makeglossaries
% ---- Acronyms & Glossary entries ----
% Define your acronyms and glossary terms here.
% Usage in text:
%   \acrfull{key}   → full form + acronym
%   \acrshort{key}  → acronym only
%   \acrlong{key}   → long form only
%   \gls{key}       → first use: full; subsequent: short

\renewcommand*{\acronymname}{الاختصارات}
\renewcommand*{\glossaryname}{المصطلحات}

% ---- Example acronyms (replace with your own) ----
\newacronym{bci}{\LR{BCI}}{\textarabic{ترابط الدماغ–الحاسوب} - \LR{Brain-Computer Interface}}
\newacronym{eeg}{\LR{EEG}}{\textarabic{تخطيط الدماغ الكهربائي} - \LR{Electroencephalography}}
\newacronym{ai}{\LR{AI}}{\textarabic{الذكاء الاصطناعي} - \LR{Artificial Intelligence}}
\newacronym{ml}{\LR{ML}}{\textarabic{تعلُّم الآلة} - \LR{Machine Learning}}

% ---- Example glossary entries ----
\newglossaryentry{neuron}{
  name=\LR{Neuron},
  description={العَصَبُون: الوحدة الوظيفية الأساسية للجهاز العصبي}
}


% ---- Captions in Arabic ----
\addto\captionsarabic{%
  \renewcommand{\contentsname}{الفهرس}%
  \renewcommand{\figurename}{الشكل}%
  \renewcommand{\tablename}{الجدول}%
  \renewcommand{\refname}{المراجع}%
  \renewcommand{\abstractname}{الملخص}%
}

% ---- Section formatting ----
\titleformat{\section}{\Large\bfseries}{\thesection}{0.7em}{}
\titleformat{\subsection}{\large\bfseries}{\thesubsection}{0.6em}{}
\titleformat{\subsubsection}{\normalsize\bfseries}{\thesubsubsection}{0.5em}{}
\titlespacing*{\section}{0pt}{1.4ex plus 0.4ex}{0.9ex plus 0.3ex}
\titlespacing*{\subsection}{0pt}{1.0ex plus 0.3ex}{0.7ex plus 0.2ex}

% ---- Headers / footers ----
\pagestyle{fancy}
\fancyhf{}
\fancyhead[R]{\small\itshape عنوان مختصر للمقال}
\fancyhead[L]{\small د. اسم المؤلف}
\fancyfoot[C]{\small\thepage}
\renewcommand{\headrulewidth}{0.3pt}
\renewcommand{\footrulewidth}{0pt}

% ---- List spacing ----
\setlist[itemize]{topsep=0.3em, itemsep=0.15em, parsep=0pt}
\setlist[enumerate]{topsep=0.3em, itemsep=0.15em, parsep=0pt}

% ---- Table column types ----
\newcolumntype{L}[1]{>{\raggedright\arraybackslash}p{#1}}
\newcolumntype{C}[1]{>{\centering\arraybackslash}p{#1}}

% ---- Line spacing ----
\setstretch{1.5}

\begin{document}

% ============================================================
% TITLE BLOCK
% ============================================================
\begin{center}
{\LARGE\bfseries عنوان المقال العلمي}\\[1em]
{\large د. اسم المؤلف}\\[0.3em]
{\small القسم، الجامعة}\\[0.3em]
{\small \texttt{email@example.com}}\\[1em]
\end{center}

\vspace{0.5em}

% ============================================================
% ABSTRACT
% ============================================================
\begin{center}
\begin{minipage}{0.85\textwidth}
\noindent\textbf{الملخص:}
% TODO: اكتب ملخص المقال هنا (150-250 كلمة). اذكر المشكلة، المنهجية، النتائج، والاستنتاج.

\vspace{0.5em}
\noindent\textbf{الكلمات المفتاحية:} كلمة 1، كلمة 2، كلمة 3، كلمة 4.
\end{minipage}
\end{center}

\vspace{1em}

% ============================================================
% BODY
% ============================================================

\section{المقدمة}
\label{sec:intro}

% TODO: اكتب مقدمة المقال هنا. استخدم \acrfull{key} و \gls{key} للإشارة إلى المصطلحات.
% استخدم \cite{bibkey} للإشارة إلى المراجع.

\section{الخلفية النظرية والأدبيات}
\label{sec:background}

% TODO: اكتب مراجعة الأدبيات هنا.

\subsection{المبحث الفرعي}

% TODO: اكتب المحتوى هنا.

\section{المنهجية}
\label{sec:method}

% TODO: اكتب منهجية البحث هنا.

\section{النتائج}
\label{sec:results}

% TODO: اكتب النتائج هنا.

% ---- Example table ----
% \begin{table}[H]
% \centering
% \caption{عنوان الجدول}
% \begin{tabularx}{\textwidth}{L{4cm} X}
% \toprule
% \textbf{العنوان} & \textbf{الوصف}\\
% \midrule
% بيانات & بيانات\\
% \bottomrule
% \end{tabularx}
% \end{table}

% ---- Example figure ----
% \begin{figure}[H]
% \centering
% \includegraphics[width=0.8\textwidth]{figure.png}
% \caption{عنوان الشكل}
% \end{figure}

\section{المناقشة}
\label{sec:discussion}

% TODO: اكتب مناقشة النتائج هنا.

\section{الاستنتاجات}
\label{sec:conclusion}

% TODO: اكتب الاستنتاجات هنا.

% ============================================================
% APPENDICES (optional)
% ============================================================
\appendix
\section*{الملحق أ}
\addcontentsline{toc}{section}{الملحق أ}

% TODO: اكتب محتوى الملحق هنا.

% ============================================================
% GLOSSARY
% ============================================================
\printglossaries

% ============================================================
% BIBLIOGRAPHY
% ============================================================
% Add your .bib file as bibliography.bib in the project folder.
% Run: xelatex → bibtex → xelatex → xelatex
\bibliography{bibliography}
\bibliographystyle{apalike}

\end{document}
